\documentclass[conference]{IEEEtran}

\usepackage{cite}
\usepackage{amsmath,amssymb,amsfonts}
\usepackage{booktabs}
\usepackage{xcolor}
\usepackage{url}
\usepackage{listings}

\lstset{
  basicstyle=\ttfamily\scriptsize,
  keywordstyle=\bfseries,
  commentstyle=\color{gray},
  frame=single,
  breaklines=true,
  showstringspaces=false,
  tabsize=2
}

\begin{document}

\title{SentinelAPI: A Zero-Trust API Vulnerability Scanner via Cross-Principal Differential Testing}

\author{
\IEEEauthorblockN{Team SentinelAPI}
\IEEEauthorblockA{\textit{AMIHACKS 2026 --- Track C (Industry / Deep-Tech)}\\
Jaipur, Rajasthan, India}
}

\maketitle

\begin{abstract}
Broken object-level authorization (BOLA/IDOR) and excessive data exposure are
the two most damaging classes of API vulnerability, yet most engineering teams
review API security manually, infrequently, and only after deployment. We
present \textbf{SentinelAPI}, an automated scanner that detects five
vulnerability classes---BOLA, excessive data exposure, broken authentication,
broken function-level authorization, and mass assignment---from an OpenAPI
specification alone. The core of the approach is \emph{differential testing
across security contexts}: the scanner registers two synthetic principals and
systematically injects one principal's identifiers into the other's session,
flagging any endpoint that returns cross-principal data. Independently, it
diffs the actual response against the schema declared in the specification to
detect over-exposure of sensitive fields, and it compares declared security
requirements against observed enforcement. The engine is implemented twice, in
Go and Python, behind a common command-line boundary, and is verified against a
corpus of two deliberately vulnerable sandbox APIs with \textbf{precision =
recall = F1 = 1.0} (24/24 findings, zero false positives, zero false negatives).
\end{abstract}

\begin{IEEEkeywords}
API security, broken object-level authorization, differential testing, OpenAPI,
vulnerability scanner
\end{IEEEkeywords}

\section{Introduction}
Modern applications are assembled from dozens to hundreds of internal and
third-party APIs. Breaches that exploit broken object-level authorization
(``API1:2023'') and excessive data exposure are a leading cause of production
data leakage \cite{owasp}. The defect is subtle: a request can be syntactically
valid and authenticated yet still return another user's data, because the flaw
lives in \emph{authorization logic}, not request formatting. Signature-based web
scanners miss this class entirely, enterprise API-security platforms are
inaccessible to smaller teams, and manual penetration testing does not scale
with deployment velocity.

SentinelAPI closes this gap with a lightweight, deterministic scanner whose
correctness is verifiable by construction: it encodes the API's \emph{security
contract} as a set of oracles and tests the contract against observed behavior,
rendering explainable, severity-ranked findings with copy-paste reproduction.

\section{Related Work}
\emph{Differential testing} \cite{mckeeman} runs the same input through two or
more contexts and treats divergence as a defect. SentinelAPI applies this to
security: the same endpoint is executed under two principals, and any
difference in data ownership is a defect.

\emph{RESTler} \cite{restler} generates stateful sequences of API calls from a
Swagger/OpenAPI spec to find reliability bugs, using a dynamically inferred
dependency graph. \emph{EvoMaster} \cite{evomaster} applies evolutionary search
to REST API test generation. \emph{Morest} \cite{morest} adds model-based state
fuzzing. These tools target reliability and coverage; our goal is narrower and
our oracle is stronger---we test a specific, well-defined authorization property
rather than search for crashes.

\section{Threat Model and Method}
\subsection{Threat model}
We assume an attacker who is an \emph{authenticated but unauthorized} user: they
possess a valid token but should not be able to read or modify other
principals' objects. This matches the most common real-world API breach profile.

\subsection{Cross-principal differential oracle (BOLA)}
The scanner registers two synthetic principals $A$ and $B$. Registration returns
each principal's identifying fields, inferred generically from the response
rather than hardcoded. For every endpoint $e$ with a path parameter $p$, the
scanner substitutes each of $B$'s identifiers $v$ into $p$ while presenting
$A$'s token, and inspects the response for $B$-owned data:
\begin{equation}
\mathrm{leak}(e,v) \;=\; \mathrm{HTTP}\,200\big(R(A,e[p{:=}v])\big)
\;\wedge\;
\exists\, m \in M_B : m \in R(A,e[p{:=}v])
\end{equation}
where $R(A,x)$ is the response to request $x$ under $A$'s session and $M_B$ is
the set of $B$'s identifying values. A satisfied predicate flags a CRITICAL
finding. The oracle inspects \emph{content}, not status code alone---a 404 is
safe behavior, whereas a 200 containing foreign data is a leak.

\subsection{Schema-contract exposure oracle}
The specification declares, for each endpoint, a response schema---the set of
fields the server \emph{promises} to return. The scanner computes the set
difference between the paths actually returned and the paths declared, recursing
into nested objects:
\begin{equation}
\mathrm{extra}(e) \;=\;
\mathrm{paths}\big(R(A, e[p{:=}\mathrm{id}_A])\big)
\;\setminus\; \mathrm{schemaPaths}(e)
\end{equation}
Any undeclared path whose leaf is sensitive (password, ssn, card, token,
\ldots) is flagged HIGH; others MEDIUM.

\subsection{Broken authentication oracle}
When the specification declares an API-key security requirement for an endpoint,
the scanner probes the endpoint with \emph{no} Authorization header. If the
endpoint returns data (HTTP 200), the declared security is not enforced---a
broken-authentication finding.

\subsection{Broken function-level authorization (BFLA)}
For endpoints under an administrative path prefix, the scanner probes with a
regular (non-admin) token. A 200 response indicates a missing role check.

\subsection{Mass assignment}
For write endpoints (POST/PUT/PATCH), the scanner injects privileged fields
(\texttt{role}, \texttt{is\_admin}) into the request body and flags any endpoint
that accepts and reflects them.

\section{System Architecture}
The pipeline is: \emph{parse OpenAPI $\rightarrow$ register principals
$\rightarrow$ generate tests $\rightarrow$ execute across security contexts
$\rightarrow$ oracle evaluation $\rightarrow$ severity-ranked report}. The engine
is isolated behind a command-line boundary (\texttt{--base} $\rightarrow$
\texttt{findings.json} + \texttt{report.html}), so orchestration is
language-agnostic. The product wraps the engine in a FastAPI service exposing a
consent-gated, SSRF-guarded scan API and a web dashboard.

\section{Implementation}
Two engines implement identical logic: a \textbf{Go} engine (stdlib-only,
single static binary) and a \textbf{Python} reference (httpx). Identifier
inference is generic---any scalar field returned by the registration endpoint
becomes a candidate identifier---so the scanner is target-agnostic. Every
finding carries severity, endpoint, a plain-language explanation, a copy-paste
\texttt{curl} reproduction, and a JSON evidence diff.

\section{Evaluation}
We evaluated against a corpus of two deliberately vulnerable sandbox APIs (a
``Bank'' API and a ``Library'' API with different identifier names and flaw
mixes), each including secure control endpoints that must produce zero findings.

\begin{table}[h]
\centering
\caption{Benchmark results (both engines)}
\begin{tabular}{lcc}
\toprule
\textbf{Metric} & \textbf{Bank} & \textbf{Library} \\
\midrule
Vulnerability classes & 5 & 3 \\
Seeded findings & 19 & 5 \\
Detected & 19 & 5 \\
False positives & 0 & 0 \\
False negatives & 0 & 0 \\
\bottomrule
\end{tabular}
\label{tab:results}
\end{table}

Aggregated over both targets, the scanner achieves
\textbf{precision = recall = F1 = 1.0} ($24/24$ true positives). The secure
control endpoints (ownership-enforced reads, admin-gated stats) produced zero
findings, confirming absence of false positives on correctly implemented
endpoints.

\section{Novel Contributions}
\begin{itemize}
\item \textbf{A spec-derived authorization oracle}---the security contract is
  encoded as a differential test, not guessed by fuzzing.
\item \textbf{Generic principal-identifier inference}---no hardcoded field names,
  enabling the scanner to generalize across API designs.
\item \textbf{Recursive schema-contract diffing}---detects nested over-exposure,
  not just top-level field leaks.
\item \textbf{Declared-vs-enforced security comparison}---detects broken
  authentication from the specification's own security declarations.
\item \textbf{Explainability by construction}---every finding ships a
  reproduction and evidence diff, yielding zero false-positive noise.
\item \textbf{A two-target benchmark harness} reporting precision/recall/F1
  against seeded ground truth.
\end{itemize}

\section{Limitations and Future Work}
The scanner currently models single-principal per-endpoint access; role lattices,
session (shadow-token) handling, and side-channel timing remain future work, as
do LLM-assisted oracles for business-logic leaks and CI/CD integration.

\section{Conclusion}
SentinelAPI demonstrates that the most damaging API vulnerability classes can be
detected deterministically, with high precision and full explainability, by
treating API security as a differential-testing problem against the API's own
specification. The working system detects five vulnerability classes with F1 =
1.0 across a two-target benchmark, with a clear path to a Rust-class
high-concurrency engine.

\begin{thebibliography}{9}
\bibitem{owasp}
OWASP Foundation. \emph{OWASP API Security Top 10:2023}. \url{https://owasp.org/API-Security/}
\bibitem{mckeeman}
W.~M. McKeeman, ``Differential Testing for Software,'' \emph{Digital Technical Journal}, vol.~10, no.~1, pp.~100--107, 1998.
\bibitem{restler}
V.~Atlidakis, P.~Godefroid, and M.~Polishchuk, ``RESTler: Stateful REST API Fuzzing,'' in \emph{Proc. IEEE/ACM ICSE}, 2019.
\bibitem{evomaster}
A.~Arcuri, ``RESTful API Automated Test Case Generation with EvoMaster,'' \emph{ACM Trans. Softw. Eng. Methodol.}, vol.~28, no.~1, 2019.
\bibitem{morest}
Y.~Liu \emph{et al.}, ``Morest: Model-based RESTful API Testing with Execution Feedback,'' in \emph{Proc. ISSTA}, 2022.
\end{thebibliography}

\end{document}
